% 数据依赖图
% 使用方法: 在主文档中 % 数据依赖图
% 使用方法: 在主文档中 % 数据依赖图
% 使用方法: 在主文档中 % 数据依赖图
% 使用方法: 在主文档中 \input{figures/data_dependency.tex}
% 字体: 中文使用宋体（继承文档设置），英文使用 Times New Roman（newtxtext）

\begin{tikzpicture}[
    every node/.append style={font=\rmfamily},  % 确保使用衬线字体（Times/宋体）
    data/.style={draw, thick, rounded corners=3pt, fill=blue!10, minimum width=1.8cm, minimum height=0.6cm, align=center, font=\footnotesize},
    proc/.style={draw, thick, rounded corners=3pt, fill=green!15, minimum width=2.2cm, minimum height=0.7cm, align=center, font=\footnotesize},
    output/.style={draw, thick, rounded corners=3pt, fill=orange!15, minimum width=1.5cm, minimum height=0.6cm, align=center, font=\footnotesize},
    final/.style={draw, thick, rounded corners=3pt, fill=red!15, minimum width=2.5cm, minimum height=0.7cm, align=center, font=\footnotesize},
    arrow/.style={->, >=stealth, thick},
    group_box/.style={draw, dashed, thick, rounded corners=5pt},
]

% ========== 输入数据 ==========
\node[data] (state_quat) at (0, 0) {状态姿态\\参数};
\node[data] (state_pos) at (0, -1) {状态位置\\参数};
\node[data] (pts_pw) at (0, -2) {路标点\\世界坐标};
\node[data] (intrin) at (0, -3) {内参\\参数};

\node[group_box, fit=(state_quat)(state_pos)(pts_pw)(intrin), inner sep=5pt] (input_box) {};
\node[anchor=south, font=\scriptsize\bfseries] at (input_box.north) {输入数据};

% ========== 坐标变换 ==========
\node[proc] (coord_trans) at (4, -1.5) {坐标变换\\$\mathbf{P}_w \to \mathbf{P}_c$};

\draw[arrow] (state_quat) -- (coord_trans);
\draw[arrow] (state_pos) -- (coord_trans);
\draw[arrow] (pts_pw) -- (coord_trans);
\draw[arrow] (intrin) -- (coord_trans);

% ========== 残差计算 ==========
\node[output] (Pc) at (7, -1) {$\mathbf{P}_c$};
\node[proc] (residual) at (10, -1) {残差计算};
\node[output] (r) at (13, -1) {$\mathbf{r}$};

\draw[arrow] (coord_trans) -- (Pc);
\draw[arrow] (Pc) -- (residual);
\draw[arrow] (residual) -- (r);

% ========== 雅可比计算 ==========
\node[proc] (jacobian) at (7, -3) {雅可比计算};

\draw[arrow] (Pc) -- (jacobian);
\draw[arrow] (coord_trans) |- (jacobian);

% ========== 雅可比输出 ==========
\node[output] (jx) at (4, -4.5) {$\mathbf{J}_p$};
\node[output] (jf) at (7, -4.5) {$\mathbf{J}_m$};
\node[output] (dr_dpc) at (10, -4.5) {$\frac{\partial r}{\partial p_c}$};

\draw[arrow] (jacobian) -- (jx);
\draw[arrow] (jacobian) -- (jf);
\draw[arrow] (jacobian) -- (dr_dpc);

% ========== Hessian 累积更新 ==========
\node[group_box, minimum width=10cm, minimum height=1.8cm, fill=yellow!5] (hessian_box) at (7, -6.5) {};
\node[anchor=north, font=\scriptsize\bfseries] at (7, -5.7) {Hessian 累积更新};

\node[output] (A) at (3.5, -6.8) {$\mathbf{H}_{pp}$\\(累加)};
\node[output] (b) at (5.5, -6.8) {$\mathbf{b}_p$\\(累加)};
\node[output] (V) at (7.5, -6.8) {$\mathbf{H}_{mm}$\\(累加)};
\node[output] (g) at (9.5, -6.8) {$\mathbf{b}_m$\\(累加)};

\draw[arrow] (jx) -- (A);
\draw[arrow] (jx) -- (b);
\draw[arrow] (r.south) -- ++(0, -1) -| (b);
\draw[arrow] (jf) -- (V);
\draw[arrow] (jf) -- (g);
\draw[arrow] (r.south) -- ++(0, -1) -| (g);

% H_{pm} 矩阵
\node[output] (W) at (11.5, -6.8) {$\mathbf{H}_{pm}$};
\draw[arrow] (jx.south) -- ++(0, -0.5) -| (W);
\draw[arrow] (jf.south) -- ++(0, -0.5) -| (W);

% ========== Schur 消元 ==========
\node[final] (schur) at (7, -8.5) {Schur消元\\等待 Hessian 累积完成};

\draw[arrow] (A) -- (schur);
\draw[arrow] (b) -- (schur);
\draw[arrow] (V) -- (schur);
\draw[arrow] (g) -- (schur);
\draw[arrow] (W) -- (schur);

% ========== 最终输出 ==========
\node[final, fill=green!20] (output) at (7, -10) {$\mathbf{S}, \mathbf{b}'_p$\\(输出)};

\draw[arrow] (schur) -- (output);

\end{tikzpicture}

% 字体: 中文使用宋体（继承文档设置），英文使用 Times New Roman（newtxtext）

\begin{tikzpicture}[
    every node/.append style={font=\rmfamily},  % 确保使用衬线字体（Times/宋体）
    data/.style={draw, thick, rounded corners=3pt, fill=blue!10, minimum width=1.8cm, minimum height=0.6cm, align=center, font=\footnotesize},
    proc/.style={draw, thick, rounded corners=3pt, fill=green!15, minimum width=2.2cm, minimum height=0.7cm, align=center, font=\footnotesize},
    output/.style={draw, thick, rounded corners=3pt, fill=orange!15, minimum width=1.5cm, minimum height=0.6cm, align=center, font=\footnotesize},
    final/.style={draw, thick, rounded corners=3pt, fill=red!15, minimum width=2.5cm, minimum height=0.7cm, align=center, font=\footnotesize},
    arrow/.style={->, >=stealth, thick},
    group_box/.style={draw, dashed, thick, rounded corners=5pt},
]

% ========== 输入数据 ==========
\node[data] (state_quat) at (0, 0) {状态姿态\\参数};
\node[data] (state_pos) at (0, -1) {状态位置\\参数};
\node[data] (pts_pw) at (0, -2) {路标点\\世界坐标};
\node[data] (intrin) at (0, -3) {内参\\参数};

\node[group_box, fit=(state_quat)(state_pos)(pts_pw)(intrin), inner sep=5pt] (input_box) {};
\node[anchor=south, font=\scriptsize\bfseries] at (input_box.north) {输入数据};

% ========== 坐标变换 ==========
\node[proc] (coord_trans) at (4, -1.5) {坐标变换\\$\mathbf{P}_w \to \mathbf{P}_c$};

\draw[arrow] (state_quat) -- (coord_trans);
\draw[arrow] (state_pos) -- (coord_trans);
\draw[arrow] (pts_pw) -- (coord_trans);
\draw[arrow] (intrin) -- (coord_trans);

% ========== 残差计算 ==========
\node[output] (Pc) at (7, -1) {$\mathbf{P}_c$};
\node[proc] (residual) at (10, -1) {残差计算};
\node[output] (r) at (13, -1) {$\mathbf{r}$};

\draw[arrow] (coord_trans) -- (Pc);
\draw[arrow] (Pc) -- (residual);
\draw[arrow] (residual) -- (r);

% ========== 雅可比计算 ==========
\node[proc] (jacobian) at (7, -3) {雅可比计算};

\draw[arrow] (Pc) -- (jacobian);
\draw[arrow] (coord_trans) |- (jacobian);

% ========== 雅可比输出 ==========
\node[output] (jx) at (4, -4.5) {$\mathbf{J}_p$};
\node[output] (jf) at (7, -4.5) {$\mathbf{J}_m$};
\node[output] (dr_dpc) at (10, -4.5) {$\frac{\partial r}{\partial p_c}$};

\draw[arrow] (jacobian) -- (jx);
\draw[arrow] (jacobian) -- (jf);
\draw[arrow] (jacobian) -- (dr_dpc);

% ========== Hessian 累积更新 ==========
\node[group_box, minimum width=10cm, minimum height=1.8cm, fill=yellow!5] (hessian_box) at (7, -6.5) {};
\node[anchor=north, font=\scriptsize\bfseries] at (7, -5.7) {Hessian 累积更新};

\node[output] (A) at (3.5, -6.8) {$\mathbf{H}_{pp}$\\(累加)};
\node[output] (b) at (5.5, -6.8) {$\mathbf{b}_p$\\(累加)};
\node[output] (V) at (7.5, -6.8) {$\mathbf{H}_{mm}$\\(累加)};
\node[output] (g) at (9.5, -6.8) {$\mathbf{b}_m$\\(累加)};

\draw[arrow] (jx) -- (A);
\draw[arrow] (jx) -- (b);
\draw[arrow] (r.south) -- ++(0, -1) -| (b);
\draw[arrow] (jf) -- (V);
\draw[arrow] (jf) -- (g);
\draw[arrow] (r.south) -- ++(0, -1) -| (g);

% H_{pm} 矩阵
\node[output] (W) at (11.5, -6.8) {$\mathbf{H}_{pm}$};
\draw[arrow] (jx.south) -- ++(0, -0.5) -| (W);
\draw[arrow] (jf.south) -- ++(0, -0.5) -| (W);

% ========== Schur 消元 ==========
\node[final] (schur) at (7, -8.5) {Schur消元\\等待 Hessian 累积完成};

\draw[arrow] (A) -- (schur);
\draw[arrow] (b) -- (schur);
\draw[arrow] (V) -- (schur);
\draw[arrow] (g) -- (schur);
\draw[arrow] (W) -- (schur);

% ========== 最终输出 ==========
\node[final, fill=green!20] (output) at (7, -10) {$\mathbf{S}, \mathbf{b}'_p$\\(输出)};

\draw[arrow] (schur) -- (output);

\end{tikzpicture}

% 字体: 中文使用宋体（继承文档设置），英文使用 Times New Roman（newtxtext）

\begin{tikzpicture}[
    every node/.append style={font=\rmfamily},  % 确保使用衬线字体（Times/宋体）
    data/.style={draw, thick, rounded corners=3pt, fill=blue!10, minimum width=1.8cm, minimum height=0.6cm, align=center, font=\footnotesize},
    proc/.style={draw, thick, rounded corners=3pt, fill=green!15, minimum width=2.2cm, minimum height=0.7cm, align=center, font=\footnotesize},
    output/.style={draw, thick, rounded corners=3pt, fill=orange!15, minimum width=1.5cm, minimum height=0.6cm, align=center, font=\footnotesize},
    final/.style={draw, thick, rounded corners=3pt, fill=red!15, minimum width=2.5cm, minimum height=0.7cm, align=center, font=\footnotesize},
    arrow/.style={->, >=stealth, thick},
    group_box/.style={draw, dashed, thick, rounded corners=5pt},
]

% ========== 输入数据 ==========
\node[data] (state_quat) at (0, 0) {状态姿态\\参数};
\node[data] (state_pos) at (0, -1) {状态位置\\参数};
\node[data] (pts_pw) at (0, -2) {路标点\\世界坐标};
\node[data] (intrin) at (0, -3) {内参\\参数};

\node[group_box, fit=(state_quat)(state_pos)(pts_pw)(intrin), inner sep=5pt] (input_box) {};
\node[anchor=south, font=\scriptsize\bfseries] at (input_box.north) {输入数据};

% ========== 坐标变换 ==========
\node[proc] (coord_trans) at (4, -1.5) {坐标变换\\$\mathbf{P}_w \to \mathbf{P}_c$};

\draw[arrow] (state_quat) -- (coord_trans);
\draw[arrow] (state_pos) -- (coord_trans);
\draw[arrow] (pts_pw) -- (coord_trans);
\draw[arrow] (intrin) -- (coord_trans);

% ========== 残差计算 ==========
\node[output] (Pc) at (7, -1) {$\mathbf{P}_c$};
\node[proc] (residual) at (10, -1) {残差计算};
\node[output] (r) at (13, -1) {$\mathbf{r}$};

\draw[arrow] (coord_trans) -- (Pc);
\draw[arrow] (Pc) -- (residual);
\draw[arrow] (residual) -- (r);

% ========== 雅可比计算 ==========
\node[proc] (jacobian) at (7, -3) {雅可比计算};

\draw[arrow] (Pc) -- (jacobian);
\draw[arrow] (coord_trans) |- (jacobian);

% ========== 雅可比输出 ==========
\node[output] (jx) at (4, -4.5) {$\mathbf{J}_p$};
\node[output] (jf) at (7, -4.5) {$\mathbf{J}_m$};
\node[output] (dr_dpc) at (10, -4.5) {$\frac{\partial r}{\partial p_c}$};

\draw[arrow] (jacobian) -- (jx);
\draw[arrow] (jacobian) -- (jf);
\draw[arrow] (jacobian) -- (dr_dpc);

% ========== Hessian 累积更新 ==========
\node[group_box, minimum width=10cm, minimum height=1.8cm, fill=yellow!5] (hessian_box) at (7, -6.5) {};
\node[anchor=north, font=\scriptsize\bfseries] at (7, -5.7) {Hessian 累积更新};

\node[output] (A) at (3.5, -6.8) {$\mathbf{H}_{pp}$\\(累加)};
\node[output] (b) at (5.5, -6.8) {$\mathbf{b}_p$\\(累加)};
\node[output] (V) at (7.5, -6.8) {$\mathbf{H}_{mm}$\\(累加)};
\node[output] (g) at (9.5, -6.8) {$\mathbf{b}_m$\\(累加)};

\draw[arrow] (jx) -- (A);
\draw[arrow] (jx) -- (b);
\draw[arrow] (r.south) -- ++(0, -1) -| (b);
\draw[arrow] (jf) -- (V);
\draw[arrow] (jf) -- (g);
\draw[arrow] (r.south) -- ++(0, -1) -| (g);

% H_{pm} 矩阵
\node[output] (W) at (11.5, -6.8) {$\mathbf{H}_{pm}$};
\draw[arrow] (jx.south) -- ++(0, -0.5) -| (W);
\draw[arrow] (jf.south) -- ++(0, -0.5) -| (W);

% ========== Schur 消元 ==========
\node[final] (schur) at (7, -8.5) {Schur消元\\等待 Hessian 累积完成};

\draw[arrow] (A) -- (schur);
\draw[arrow] (b) -- (schur);
\draw[arrow] (V) -- (schur);
\draw[arrow] (g) -- (schur);
\draw[arrow] (W) -- (schur);

% ========== 最终输出 ==========
\node[final, fill=green!20] (output) at (7, -10) {$\mathbf{S}, \mathbf{b}'_p$\\(输出)};

\draw[arrow] (schur) -- (output);

\end{tikzpicture}

% 字体: 中文使用宋体（继承文档设置），英文使用 Times New Roman（newtxtext）

\begin{tikzpicture}[
    every node/.append style={font=\rmfamily},  % 确保使用衬线字体（Times/宋体）
    data/.style={draw, thick, rounded corners=3pt, fill=blue!10, minimum width=1.8cm, minimum height=0.6cm, align=center, font=\footnotesize},
    proc/.style={draw, thick, rounded corners=3pt, fill=green!15, minimum width=2.2cm, minimum height=0.7cm, align=center, font=\footnotesize},
    output/.style={draw, thick, rounded corners=3pt, fill=orange!15, minimum width=1.5cm, minimum height=0.6cm, align=center, font=\footnotesize},
    final/.style={draw, thick, rounded corners=3pt, fill=red!15, minimum width=2.5cm, minimum height=0.7cm, align=center, font=\footnotesize},
    arrow/.style={->, >=stealth, thick},
    group_box/.style={draw, dashed, thick, rounded corners=5pt},
]

% ========== 输入数据 ==========
\node[data] (state_quat) at (0, 0) {状态姿态\\参数};
\node[data] (state_pos) at (0, -1) {状态位置\\参数};
\node[data] (pts_pw) at (0, -2) {路标点\\世界坐标};
\node[data] (intrin) at (0, -3) {内参\\参数};

\node[group_box, fit=(state_quat)(state_pos)(pts_pw)(intrin), inner sep=5pt] (input_box) {};
\node[anchor=south, font=\scriptsize\bfseries] at (input_box.north) {输入数据};

% ========== 坐标变换 ==========
\node[proc] (coord_trans) at (4, -1.5) {坐标变换\\$\mathbf{P}_w \to \mathbf{P}_c$};

\draw[arrow] (state_quat) -- (coord_trans);
\draw[arrow] (state_pos) -- (coord_trans);
\draw[arrow] (pts_pw) -- (coord_trans);
\draw[arrow] (intrin) -- (coord_trans);

% ========== 残差计算 ==========
\node[output] (Pc) at (7, -1) {$\mathbf{P}_c$};
\node[proc] (residual) at (10, -1) {残差计算};
\node[output] (r) at (13, -1) {$\mathbf{r}$};

\draw[arrow] (coord_trans) -- (Pc);
\draw[arrow] (Pc) -- (residual);
\draw[arrow] (residual) -- (r);

% ========== 雅可比计算 ==========
\node[proc] (jacobian) at (7, -3) {雅可比计算};

\draw[arrow] (Pc) -- (jacobian);
\draw[arrow] (coord_trans) |- (jacobian);

% ========== 雅可比输出 ==========
\node[output] (jx) at (4, -4.5) {$\mathbf{J}_p$};
\node[output] (jf) at (7, -4.5) {$\mathbf{J}_m$};
\node[output] (dr_dpc) at (10, -4.5) {$\frac{\partial r}{\partial p_c}$};

\draw[arrow] (jacobian) -- (jx);
\draw[arrow] (jacobian) -- (jf);
\draw[arrow] (jacobian) -- (dr_dpc);

% ========== Hessian 累积更新 ==========
\node[group_box, minimum width=10cm, minimum height=1.8cm, fill=yellow!5] (hessian_box) at (7, -6.5) {};
\node[anchor=north, font=\scriptsize\bfseries] at (7, -5.7) {Hessian 累积更新};

\node[output] (A) at (3.5, -6.8) {$\mathbf{H}_{pp}$\\(累加)};
\node[output] (b) at (5.5, -6.8) {$\mathbf{b}_p$\\(累加)};
\node[output] (V) at (7.5, -6.8) {$\mathbf{H}_{mm}$\\(累加)};
\node[output] (g) at (9.5, -6.8) {$\mathbf{b}_m$\\(累加)};

\draw[arrow] (jx) -- (A);
\draw[arrow] (jx) -- (b);
\draw[arrow] (r.south) -- ++(0, -1) -| (b);
\draw[arrow] (jf) -- (V);
\draw[arrow] (jf) -- (g);
\draw[arrow] (r.south) -- ++(0, -1) -| (g);

% H_{pm} 矩阵
\node[output] (W) at (11.5, -6.8) {$\mathbf{H}_{pm}$};
\draw[arrow] (jx.south) -- ++(0, -0.5) -| (W);
\draw[arrow] (jf.south) -- ++(0, -0.5) -| (W);

% ========== Schur 消元 ==========
\node[final] (schur) at (7, -8.5) {Schur消元\\等待 Hessian 累积完成};

\draw[arrow] (A) -- (schur);
\draw[arrow] (b) -- (schur);
\draw[arrow] (V) -- (schur);
\draw[arrow] (g) -- (schur);
\draw[arrow] (W) -- (schur);

% ========== 最终输出 ==========
\node[final, fill=green!20] (output) at (7, -10) {$\mathbf{S}, \mathbf{b}'_p$\\(输出)};

\draw[arrow] (schur) -- (output);

\end{tikzpicture}
